% =====================================================================
%  R. Nielsen, LOGIK OG PSYCHOLOGIE. CURSUS FOR UNIVERSITETSAARET 1871-72.
%  Kjøbenhavn: J. H. Schubothes Boghandel (Louis Kleins Bogtrykkeri), [1871].  8 pp.  Fraktur.
%  An eight-page print for the lecture course: "Psychologie. Indledning." -- the relation of soul and body
%  (Dualisme, Identitetslære ...).  It treats the subject of the 1867-68 course's pp. 424 ff.
%  (texts/nielsen/logik-og-psychologie) in a new text.
%
%  TRANSCRIPTION -- GENERATED; DO NOT EDIT BY HAND.  Rebuild, chained (&&), from the repo root:
%    python3 ebook-method/logik-og-psychologie-1871/draft.py --emit --force
%      this template + tesseract's script/Fraktur reading of KB's scan (the scan has NO text layer; the reading is
%      set as one by ocr.py MODEL=layer); ~/bibliotek/Nielsen, Rasmus/1871-logik-og-psychologie-omtrykte-blade.pdf,
%      PDF 250-259 (PDF 2-247 of that scan are another book, Funke's Psychologie of 1799), see SCANS.tsv;
%    python3 ebook-method/collate/check.py logik-og-psychologie-1871 --witness ocr --apply
%      the decisions in decisions.tsv, taken against an independent reading by a different model (tesseract deu_latf);
%    python3 ebook-method/logik-og-psychologie-1871/oepass.py --apply
%      the ø pass and the word fixes (word-fixes.tsv);
%    python3 ebook-method/logik-og-psychologie-1871/finish.py --apply
%      the footnotes, corrections across markup, Latin (\textit), letterspacing (\emph), PRINTED AS IS marks.
%
%  WORD-ACCURACY (2026-10-13).  No e-book, no earlier transcription and no text layer exist; the 1867-68 course
%  treats the subject in a different text, so there is no finding aid.  The draft is tesseract's script/Fraktur reading
%  (tessdata_best), set as a layer; it was collated word by word, with punctuation, quote kinds and word division,
%  against tesseract's deu_latf model: 241 disagreements in 1908 words.  158 were settled by rule (rules.py, each
%  decision named: each model's known Fraktur faults undone, accepted only where the result is a word attested in
%  other books of the corpus), 10 Sjæl- forms by the third reading (tesseract dan, which alone sees æ), the rest
%  read at the image: 72 crops read BLIND by one reader (controls 2 of 2), decided by script (readsbox.py) or by the
%  main session from the reader's box and line.  OPEN 0.  Ø: oepass.py restored it where the o-form is no word (34).
%  Then a FULL blind page reading of all eight pages and the title (one reader, 1924 words; .parts/pagecheck/), compared
%  word by word with the collated text: 29 word and punctuation faults that both machine readings shared (Fraktur k
%  read as l, f, st or t; ø read as o: Folge, metaphysist, iffe, Produft ...) were set from it (finish.py POST), the
%  doubtful ones (bære, Leibnitz, the points at pp. 1 and 3, „o. s. v.“ p. 2) on line crops at 300 dpi; the headings,
%  paragraphs, antiqua and letterspacing were taken from it and checked on line crops.  After that the reading and the
%  text differ in 21 words and 38 marks, every one settled for the text (reader's slips: uensartede for ueensartede x3,
%  nedlaat, Enhed, Menneskets for „er Mennesket Aand“ p. 3 = man is mind; paragraph and heading signs), and two short letterspaced words the reader left plain (Aand, Lys).
%  NOT ESTABLISHED: misprints that both machines and the reader read alike as real words.
%  PAGE MARKERS: 8 of 8 (p. 1 carries no folio).
%
%  CONVENTIONS (as the other Fraktur books by Nielsen).  Body in Fraktur; aa (not å); ø as printed, and by the ø rule
%  (CLAUDE.md) where the o-form is no word.  Emphasis is letterspacing -> \emph{}.  Latin and French set in antiqua
%  -> \textit{}.  Fraktur I/J are one sort: I for I-words (Idee, Ikke), J for J-words (Jeg).  The double hyphen at a
%  line end is resolved; inside a compound it is "-".  "% --- p. N ---" + \opage{N} at the first word of printed
%  page N (p. 1 carries no folio).  Footnotes *), **) per page, as printed.  Folios, signature marks and the
%  library's stamps are not transcribed.  Print governs; misprints are kept and marked PRINTED AS IS.
% =====================================================================
\documentclass[12pt,a4paper,openany]{book}
\usepackage[utf8]{inputenc}
\usepackage[T1]{fontenc}
\usepackage{libertinus}
\usepackage{libertinust1math}
\usepackage{textalpha}
\usepackage{graphicx}
\DeclareUnicodeCharacter{0254}{\reflectbox{c}}  % ɔ (U+0254) = reversed c
\usepackage[danish]{babel}
\usepackage[protrusion=true,expansion=true]{microtype}
\usepackage{setspace}
\usepackage[margin=1.3in]{geometry}
\usepackage{fancyhdr}
\usepackage{hyperref}
\hypersetup{pdftitle={Logik og Psychologie. Cursus for Universitetsaaret 1871--72}, pdfauthor={Rasmus Nielsen}, hidelinks}
\setlength{\headheight}{15pt}
\pagestyle{fancy}
\fancyhf{}
\fancyhead[LE,RO]{\thepage}
\fancyhead[C]{\textit{R. Nielsen: Logik og Psychologie. Cursus 1871--72}}
\renewcommand{\thefootnote}{\ifcase\value{footnote}\or *)\or **)\or ***)\else
  \arabic{footnote})\fi}
\setstretch{1.3}
\usepackage{marginnote}
\newcommand{\opage}[1]{\marginnote{\footnotesize\textit{[#1]}}}
\newcommand{\apage}[1]{\marginnote{\footnotesize\textit{[#1]}}}

\begin{document}

% --- the title page (PDF 250), read at the image (page reader, 2026-10-13).  Three small ornaments are drawn as rules;
% KB's stamp is not transcribed.  No year is printed: KB's record gives 1871.
\begin{titlepage}
\begin{center}
\vspace*{6em}
\rule{3em}{0.4pt}\\[1.5em]
{\LARGE Logik og Psychologie.}\\[1.5em]
\rule{3em}{0.4pt}\\[1.5em]
{\large\textbf{Cursus for Universitetsaaret 1871-72}}\\[1.5em]
{\small af}\\[1em]
{\Large R. Nielsen.}\\[2em]
\rule{3em}{0.4pt}\\[2em]
\textbf{Kjøbenhavn.}\\[0.5em]
{\small Forlagt af J. H. Schubothes Boghandel.}\\[\fill]
{\small Louis Kleins Bogtrykkeri.}
\end{center}
\end{titlepage}


% ===== TEXT (pp. 1--8) =====
% --- p. 1 ---
\setcounter{footnote}{0}%
\opage{1}\begin{center}
{\Large\textbf{Psychologie.}}\\[0.4em]
\rule{4em}{0.4pt}\\[0.4em]
Indledning.
\end{center}

Hvad det i Virkeligheden betyder, at Selvhedens Princip
maa indeholde Principet for Andet, at Selvmagt og Magt
i Andet maae bekræfte hinanden, at Aanden selv maa have
Naturen til Grund for at articulere og bære de virkelige
Naturanlæg, bliver især indlysende, naar vi betragte

\begin{center}
\textbf{Forholdet mellem Sjæl og Legeme.}\\[0.6em]
α) \textbf{Dualisme.}
\end{center}

Idet Cartesius gik ud fra den Forudsætning, at Sjæl og
Legeme vare to ueensartede Substanser, Sjælen den tænkende,
Legemet den udstrakte Substans, satte han med det Samme
en saadan Skilsmisse mellem Selvhedens og Andethedens
Væsen, at Vexelvirkningen blev utænkelig. Forestillingen om
en gjensidig physisk Indvirkning (\textit{systema influxus physici})
var kun en foreløbig populær Antagelse, der snart maatte
give Plads for metaphysiske Theorier.

Hos Malebranche søges Dualismen hævet derved, at
alle Tingsaarsager nedsættes til at være Anledningsaarsager
(\textit{causæ occasionales}) for Gud som Tilværelsens eneste virkende
Aarsag (\textit{causa efficiens}). Forholdet mellem Sjæl og
Legeme bliver da som Følge heraf opfattet saaledes, at Gud
af enhver Bevægelse i Legemet tager Anledning til at bevæge
Sjælen og omvendt (\textit{systema causarum occasionalium}).
Det almægtige Selv fik ved uophørlig Indgriben i Tingenes
Orden en saa enevældig Indflydelse paa den hele Skabning,
at Naturaarsagerne forsvandt og dermed Naturen.

% --- p. 2 ---
\setcounter{footnote}{0}%
\opage{2}For at sætte Mirakeltheoriens Overgreb en Grændse
opstillede Leibnitz Læren om den forudbestemte Harmonie
(\textit{Harmonia præstabilita}). Grundtanken er, at Gud fra
Evighed af har beregnet alle Bevægelser i det Legemlige og
Bevægelserne i Sjælen saa nøiagtig, at de to iøvrigt ueensartede
Rækker Punkt for Punkt svare til hinanden indbyrdes.

Hos Leibnitz bliver Selvhedsprincipet paa eiendommelig
Maade udviklet i Læren om Monaderne. Alle tilværende
Substanser ere Krafteenheder, hver med sit Anlæg til Selvudvikling.
Sjælen er altsaa en Monade, der ifølge sit oprindelige
Anlæg er fra sit Synspunkt (\textit{suivant son point
de vue}) istand til at forestille den hele Verden (Modsætning
mellem Mikrokosme og Makrokosme).

Men hvad er Materien? Ved Besvarelsen af dette
Spørgsmaal taler Leibnitz et dobbelt Sprog, idet han snart
stiller sig paa Metaphysikens, snart paa Physikens Standpunkt.


I metaphysisk Forstand har Materien ingen Realitet;
den er fra den subjective Side kun et Phænomen af Sjælemonadens
dunkle forvirrede Forestillinger, fra den objective
Side kun et Complex af de laveste Monader (\textit{Monades nudæ}).
Fra dette Synspunkt er altsaa hele Verden et Produkt af
lutter Selvhedsprinciper.

Betragtet fra Physikens Standpunkt har Materien derimod
en umiddelbar Realitet; Andethedsprincipet kommer forsaavidt
til sin Ret, som Leibnitz ved Bestemmelsen af de mechaniske
Bevægelseslove forudsætter Tryk, Stød, Inertie o. s. v.

Dualismens yderste Consequens er, at de to Systemer
af Love, Selvhedslovene (Sjælens Morallove) og Andethedslovene
(de mechaniske Love for Bevægelsen i Legemet) ikke
lade sig forene. Den Forestilling, at Sjælen skulde forholde
sig til sit Legeme, som en Herre, der havde en Automat til
Tjener, er saa monstrøs, at den uvilkaarlig fordrer Dualismens
Opløsning.

% --- p. 3 ---
\setcounter{footnote}{0}%
\opage{3}\begin{center}β) \textbf{Identitetslære.}\end{center}

Grundlaget er givet af Spinoza. Den relative Selvstændighed,
Cartesius havde tillagt Substanserne, Sjæl og
Legeme, er i Forhold til Gud, den absolute Substans, forsvindende.
Idet de relative Substanser tabe deres Selvstændighed,
tabe de med det Samme deres Substantialitet
og forvandle sig til Modificationer (\textit{modi}) af een eneste
uendelig Substans.

Det er er den ene uendelige Substans med sine til Aand % PRINTED AS IS: er er (doubled), page reader and the OCRs
og Legeme svarende Attributer: Tænkning og Udstrækning,
der er sin egen (\textit{causa sui}) og alle Tings Aarsag.
Den ene og samme Tingenes Orden udgjør altsaa en Række
af tænkende og en Række af udstrakte Existenser (\textit{res}), hine
dannende Aandeverdenen, denne Legemverdenen.

Som Modus af Tænkningens Attribut er Mennesket
\emph{Aand}, som Modus af Udstrækningens Attribut \emph{Legeme}.
Det, hvorved Mennesket adskiller sig fra andre Ting, er med
Hensyn til Udstrækningen et mere sammensat Legeme, og som
Følge heraf, hvad Tænkningen angaaer, en mere sammensat
Forestilling. Aanden er Legemets Idee (\textit{idea corporis}). Dersom
det menneskelige Legeme var en Cirkel, saa vilde den
menneskelige Aand være hverken mere eller mindre end Ideen
af denne Cirkel.

Identitetslæren sammensmelter Selvhedens og Andethedens
Principer til et eneste idealt-realt Princip, og maa, for
at affinde sig med de virkelige Modsætninger, tage sin Tilflugt
til formelle Omskrivninger og tomme Ordforklaringer.

Opfattes Principet umiddelbart fra den reale Side, saa
have vi Materialismen. Det organiske Legeme er et Produkt
af mechanisk-chemiske Aarsagers eiendommelige Vexelvirkninger,
og det sjælelige Liv, en Modification af disse
Vexelvirkninger, at forklare ved en Omskrivning af det Legemlige
(\textit{système de la nature; l’homme machine.}) Op%
% --- p. 4 ---
\setcounter{footnote}{0}%
\opage{4}fattes det ideal-reale Princip fra det ideale Synspunkt, saa
bliver Identitetslæren speculativ.

Den skarpe dialektiske Antithese af \emph{Jeg} og \emph{Ikke-Jeg}
(Selvhed og Andethed), der giver Fichtes Videnskabslære
sin sande videnskabelige Betydning, blev i Schellings Naturphilosophie
i den Grad udvandet og neutraliseret, at Aandens
Selvstændighed gik tabt, og Naturkategorierne forvandledes til
hule Symboler i Læren om „\emph{den absolute Identitet}“.

Det er Aleenhedens uendelige Virksomhed, der frembringer
alt Liv, det naturlige saavel som det sædelige, derved
at den i en ordnet Trinrække fuldbyrder „\emph{sin historiske
Verdensproces}“ og med den sin „\emph{Selvbefrielse}“.

Paa det første (laveste) Trin virker Principet som
„\emph{Lys}“, i negativ Form som „\emph{Tyngde}“; Produktet er „\emph{den
uorganiske Verden}“.

Paa det andet Trin virker Principet som „\emph{irritabel
Drift}“, i negativ Form som „\emph{Sensation}“, Produktet er
„\emph{den organiske Natur}“.

Paa det tredie Trin virker Principet som „\emph{praktisk
Fornuftlov}“ eller som „\emph{Religion og Stat}“, i negativ
Form som „\emph{bevidst Driftlov}“; Produktet er „\emph{Menneskehedens
Liv}“.

At slige speculative Abstractioner, hvormange Omskrivninger
man end anvender, og hvorofte man end varierer det
forslidte Thema om Naturens alorganiserende Væsenhed, om
Universet, der efter mange dunkle Forberedelser og nødvendige
Adspredelser i det Elementære omsider samler sig til organisk
Eenhed og oplader sig til Selvbevidsthed, ere ufrugtbare og
i sig selv uvidenskabelige, vil fra Videnskabens nuværende
Standpunkt være iøinefaldende.

\begin{center}γ) \textbf{Identitetsdualismen.}\end{center}

Det Sande i Dualismen er, at Sjæl og Legeme have
deres Oprindelse fra modsatte, ueensartede Principer, Sjælen
% --- p. 5 ---
\setcounter{footnote}{0}%
\opage{5}fra Selvhedens, Legemet fra Andethedens Princip. At Materien
i sit Væsen er selvløs, har Physiken een Gang for
alle slaaet fast i Læren om Inertien: Inertie og Selvløshed
ere eenstydige Kategorier. Det er da ligesaa umuligt at
gjøre Livet til en Resultant af physiske Kræfters Samvirken
som at aflede Selvet af det Selvløse.

Det Sande i Identitetslæren er, at Sjæl og Legeme,
trods deres principielle Forskjellighed, i Ideen selv have
deres Oprindelse fra eet og samme Grundprincip, \emph{det guddommelige Magtprincip}.

Med fuld Ret have Dualismens Repræsentanter fremhævet,
at Modsætningen af Sjæl og Legeme nødvendig
maatte forudsætte en guddommelig Magteenhed; men da
Magtens Begreb i dets oprindelige Selvfordobling endnu
ikke var udviklet, ansaae man uden videre Magteenheden for
et umiddelbart Udtryk af den overnaturlige Skabermagt;
heri bestaaer Dualismens Grundvildfarelse.

I Magtens Rige, i Tingenes Verden, er Anlægseenheden,
den gjennem virkende Naturaarsager sig articulerende Ideemagt,
en Modsætningseenhed af ideel Formaaen og reel Bestaaen,
altsaa en Modsætningseenhed af Livsprincip og Materie.

At Livsprincipet i al sin Idealitet er et Magtprincip,
fremlyser deraf, at det indvirker paa Materien: thi et abstract
Tankevæsen kan umulig indvirke paa nogetsomhelst Materielt.
At Materien fra sin Side kan paavirkes af Livsprincipet
uden at Stofferne derfor opgive deres physisk-chemiske Natur,
beviser, at dens objective Sandhed maa søges i det Magtbegreb,
der i de physiske Egenskaber har nedlagt Muligheden
til organisk Virksomhed.

Dualismen beroer paa Modsætningen af Magtbegrebets
Fordobling gjennem Selvhedens og Andethedens Form; Identiteten
derimod paa Begrebseenhedens Realitet i begge Former.
Eenheden af Selvhedsmagten i det Sjælelige og
Andethedens Magt i det Legemlige gjør ikke de to Principer
% --- p. 6 ---
\setcounter{footnote}{0}%
\opage{6}til eet Princip; men den afgiver det Medium, hvori Livsvirksomheden
bliver reel, og Materien med sine Muligheder
ideel, det Grundmedium, uden hvilket en organisk Vexelvirkning mellem Sjæl og Legeme vilde være utænkelig.
(Physiologie og Psychologie).

\begin{center}
\rule{4em}{0.4pt}\\[0.6em]
\textbf{Psychologien i Grundtræk.}
\end{center}

Det er Bevidsthedslivet med sin individuelle Naturgrund
og organisk betingede Functioner, hvis Overgang fra Sandsebevidsthed
til Selvbevidsthed og aandeligt Liv er Psychologiens
Gjenstand. Vi inddele med Hensyn hertil Psychologien i
Læren om det selvbevidste Livs naturlige Grundyttringer, om
dets absolute Indhold og dets absolute Gyldighed.

\begin{center}
\textit{a)} \textbf{Det selvbevidste Livs naturlige Grundyttringer:}\\
\textbf{Sjæleevnerne.}
\end{center}

I skarp Modsætning til den registeragtige Opregnen af
Sjæleevner, hvormed man saa tidt overlæsser de psychologiske
Lærebøger, have nogle Psychologer (Herbart) ligefrem negtet
Begrebet Sjæleevne Realitet. Men i Princip og Anlæg er
Sjælen lutter Evne. Opgaven er at paavise de forskjellige
Former og Retninger, hvori den ene og samme Grundevne
naturligen individualiserer og forgrener sig.

Selvhedens alleralmindeligste Yttringsform er \emph{Fornemmelse}:
den fornemmer sig selv, idet den fornemmer sit Andet,
og den bestemmer sit Andet ved at fornemme sig selv. Fornemmelsen
er i passiv Umiddelbarhed en Sandsning, i
ophævet Umiddelbarhed en Reflexion, i activ Umiddelbarhed
en Selvbestemmen (Dialektik af det Active og det Passive).

Som Sandsning er den bestemte Selvfornemmelse en
\emph{Følelse} (Læren om Følelserne), Fornemmelsen af Andet en
umiddelbar Objectivering gjennem forskjelligartede \emph{Sandser}
(Syn, Hørelse o. s. v.); som Eenhed af Selvfornemmelse og
Fornemmelse af Andet omsætter den sin Umiddelbarhed i
% --- p. 7 ---
\setcounter{footnote}{0}%
\opage{7}\emph{Hukommelse} og \emph{Phantasie}, og hermed begynder Reflexionen.


I Reflexionen frigjøres Sandsebillederne ved Ordet
(Sprogets Betydning) til Forestillinger og Almeenbegreber
(Forstandslivet) og forklares til Eenhed af det Subjective og
det Objective i Tænkningen af det Almeengyldige (Fornuft;
Forhold imellem Forstand og Fornuft). Reflexionen fuldender
sig i \emph{Tænkningen}. Men allerede den Maade, hvorpaa
Vexelbestemmelsen mellem Selvfornemmelse og Fornemmelse
af Andet, mellem Selvbestemmelse og Bestemmelse af Andet,
forklares i Reflecteren og Tænken, viser, at Fornemmelsens
passive Umiddelbarhed oprindelig er activ, at med andre Ord
Selvfornemmelse nødvendig forudsætter Selvbestemmelse.

Oprindelig er Selvbestemmelsen, den naturlige Selvyttring,
at bestemme som Drift, Vilkaarligheden fremgaaer
af uvilkaarlig Tilskyndelse. Idet Driften peger mod det Objective,
sætter den Selvet i Bevægelse; Driften er Villiens
Grund. Men Selvbestemmelsen er reflecteret og klarer sig
med Reflexionen til Villie: Villien, der grunder i Driften,
ophæver igjen Driften og grunder i sig selv (Egoisme, Egoitet).
Den reflecterede Drift er \emph{Tilbøielighed}, og den hæftige
Tilbøielighed \emph{Lidenskab}. Forholdet mellem Drift og Villie
sees i Vexelbrydning af Lidenskab og Villie (Heroisme, Asketik).

\begin{center}
\textit{b)} \textbf{Det selvbevidste Livs absolute Indhold: Ideernes}\\
\textbf{Psychologie.}
\end{center}

Da Bevidstheden vækkes og udvikles i Tiden ved sandselige
Indtryk, endelige Fordringer og vexlende Interesser,
er det naturligt, at den fra Begyndelsen af maa fyldes med
endeligt Indhold. Som Selvbevidsthed har den ikke desto
mindre et oprindeligt Selv og just derfor et oprindeligt Krav
paa et selvgyldigt, absolut Indhold: kun Ideen kan som Idee
mætte den hungrige Subjectivitet (Aanden og det Aandelige).

Hvor naturlig den opvaagnende Selvbevidsthed trænger
% --- p. 8 ---
\setcounter{footnote}{0}%
\opage{8}til Ideens Indhold, sees af den Oprindelighed, hvormed
ethvert begavet Naturfolk begynder med religiøse Forbilleder
(Mythologie), sædelige Grunderkjendelser, Stræben efter at
give det Oversandselige sandselig Form og bestemme det
Almeengyldige. Det absolute Indhold antager nemlig en
grundforskjellig Form, eftersom det oprindelig opfattes af
Phantasien, af Tænkningen eller af Villien. Form og Indhold
gjennemtrænge hinanden saaledes, at den ene i sig absolute
Idee formedelst de tre Grundformer bliver til tre
Grundideer: Skjønhed, Sandhed, Godhed.

\begin{center}
\textit{c)} \textbf{Det selvbevidste Livs absolute Gyldighed: Sjælens}\\
\textbf{Udødelighed.}
\end{center}

Som Organ for evige Ideer maa Selvbevidstheden nødvendig
have det Evige i sit eget Væsen; herpaa beroer dens
absolute Gyldighed, herpaa grunder sig Læren om Sjælens
Udødelighed d.v.s. om Personlighedens evige Liv. At Sjælen
er udødelig i metaphysisk Forstand, kan allerede sees af Plante- og
Dyrelivet; men Forvisningen om en existentiel, for Menneskets
Vedkommende personlig Udødelighed kan først opnaaes
ved klar Indsigt i Forholdet mellem Ideen, navnlig det Godes
Idee, og den menneskelige Selvbevidsthed. Forsaavidt dette
Forhold i sin Almindelighed tydeliggjøres i Videnskaben, kan
man jo vistnok sige, at Sjælens Udødelighed lader sig videnskabelig
bevise; men da Virkelighedsbetingelserne for et tilkommende
Liv falde udenfor al Erfaring, og det rene Tankebeviis
følgelig ikke rækker ud over det Metaphysiske, indskrænker
Beviset sig, strengt taget, til at godtgjøre Fordringens ethiske
Gyldighed, ikke dens Opfyldelse. Den fuldkomne Forvisning
om Sjælens Udødelighed er af religiøs ethisk Natur og
tjener just til at vise Forskjellen imellem Tro og Viden,
imellem praktisk og theoretisk Erkjendelse (psychologisk nødvendig
Adskillelse af Religion og Videnskab).

\end{document}
